\section{Stable-homotopy objects}
\label{chap:stable-background}

This section defines the stable-homotopy language required by the paper.  Located May, Moss, and Toda inputs are stated later in the external-results chapter rather than treated as definitions.

\subsection{Spectra, triangles, and homology}

% SOURCE: PROJECT_BOUNDARY.md, Confirmed design decision 1; MainPaper/main.tex:253-263, 670-758.
\begin{definition}[Stable homotopy context]
  \label{def:stable-context}
  \uses{def:mathlib-shift,def:mathlib-triangle,def:mathlib-pretriangulated,def:mathlib-monoidal-category,def:mathlib-symmetric-category}
  \notready
  A stable homotopy context consists of a pointed additive homotopy category
  of spectra, suspension equivalences $\Sigma^n$, a symmetric monoidal smash
  product with sphere unit $S^0$, cofiber sequences, and distinguished
  triangles.  It includes rotation, functoriality, smash compatibility, and
  the sign and suspension conventions needed below.  Stable homotopy groups
  and their exact sequences are exposed separately so that later proofs do not
  hide them inside this context record.
\end{definition}

% SOURCE: MainPaper/main.tex:126-138 and 253-263; suspension conventions used throughout.
\begin{definition}[Sphere spectrum and suspension]
  \label{def:stable-sphere-suspension}
  \uses{def:stable-context}
  \notready
  The monoidal unit is the sphere spectrum $S^0$.  For every $n\in\mathbb Z$,
  suspension gives an equivalence $\Sigma^n$ with coherent isomorphisms
  $\Sigma^0X\simeq X$ and $\Sigma^{m+n}X\simeq\Sigma^m\Sigma^nX$, compatible
  with smash products and mapping groups.
\end{definition}

% SOURCE: MainPaper/main.tex:102-170 and throughout Sections 2--7.
\begin{definition}[Stable homotopy groups]
  \label{def:stable-homotopy-group}
  \uses{def:stable-sphere-suspension}
  \notready
  For a spectrum $X$ and $n\in\mathbb Z$, define
  \[
    \pi_nX=[\Sigma^nS^0,X].
  \]
  These are abelian groups, functorial in $X$, with suspension and composition
  isomorphisms fixed by the context's sign convention.
\end{definition}

% SOURCE: MainPaper/main.tex:253-263, global standing assumptions.
\begin{definition}[Admissible spectrum]
  \label{def:admissible-spectrum}
  \uses{def:stable-context}
  \notready
  A spectrum is admissible when it is $2$-complete, connective, and of finite
  type.  Strong convergence is deliberately not part of this predicate:
  whenever a classical or synthetic Adams spectral sequence is invoked, its
  convergence witness is supplied separately.  This avoids defining
  admissibility in terms of spectral sequences that are themselves indexed by
  admissible spectra.
\end{definition}

% SOURCE: MainPaper/main.tex:135-159, where products are used for the sphere;
% standard associative-algebra-object structure in a symmetric monoidal category.
\begin{definition}[Unital multiplication data]
  \label{def:unital-multiplication-data}
  \uses{def:stable-context,def:stable-sphere-suspension}
  \notready
  A unital multiplication datum on a spectrum $R$ is an associative algebra
  object in the stable symmetric monoidal context: it includes multiplication
  $R\wedge R\to R$, a unit $S^0\to R$, and all associativity and unit
  coherences.  This is extra structure and is not implied by admissibility.
  The sphere $S^0$, as the monoidal unit, has its canonical such datum.  The
  same terminology will be used for associative algebra objects in the
  synthetic context.
\end{definition}

% SOURCE: MainPaper/main.tex:688-769; Pstragowski interface used in Proposition 3.1.
\begin{definition}[$H\mathbb F_2$ homology]
  \label{def:hf2-homology}
  \uses{def:stable-context,def:mathlib-f2}
  \notready
  The context contains an Eilenberg--Mac Lane ring spectrum $H\mathbb F_2$ and
  a homology functor $H_*(-;\mathbb F_2)$ from spectra to graded
  $\mathbb F_2$-vector spaces.  It preserves suspension and sends cofiber
  sequences to long exact sequences.
\end{definition}

% SOURCE: MainPaper/main.tex:929-964 and all uses of distinguished triangles.
\begin{definition}[Cofiber and distinguished triangle data]
  \label{def:cofiber-triangle}
  \uses{def:stable-context}
  \notready
  For every map $f:X\to Y$, choose a cofiber $C(f)$ and a distinguished
  triangle
  \[
    X\xrightarrow fY\longrightarrow C(f)\longrightarrow\Sigma X.
  \]
  A morphism of triangles is three compatible maps, and a triangle equivalence
  is a morphism whose three components are isomorphisms.  Rotations and
  suspensions preserve distinguished triangles.  Every distinguished triangle
  whose first map is $f$ also carries a chosen triangle equivalence to this
  selected cofiber triangle; transporting along that equivalence is explicit
  data rather than definitional equality of its third object with $C(f)$.
\end{definition}

\begin{proposition}[Extension along a shifted cofiber inclusion]
  \label{prop:shifted-cofiber-extension}
  \lean{KIP126.StableHomotopy.CofiberExtension.shiftedCofiberTriangle,KIP126.StableHomotopy.CofiberExtension.shiftedCofiberTriangle_distinguished,KIP126.StableHomotopy.CofiberExtension.exists_extension_of_shift_comp_zero,KIP126.StableHomotopy.CofiberExtension.exists_eta_nu_extension}
  \uses{def:cofiber-triangle,def:stable-sphere-suspension,def:mathlib-pretriangulated}
  \leanok
  In any supplied stable-homotopy context with the chosen cofibers, let
  $f:X\to Y$, $n\in\mathbb Z$, and $a:\Sigma^nY\to Z$. If
  $a\circ\Sigma^nf=0$, there exists $g:\Sigma^nC(f)\to Z$ with
  $g\circ\Sigma^n\iota_f=a$. All integers and all objects are retained.
  In particular, for the same maps $\eta:S^1\to S^0$ and
  $\nu:S^3\to S^0$, the explicit premise
  $\nu\circ\Sigma^3\eta=0$ gives an extension
  $\Sigma^3C\eta\to S^0$ restricting to $\nu$.
  This does not prove the composite-zero premise for the fixed model,
  choose an actual CW comparison, or identify native coordinates.
\end{proposition}
\begin{proof}
  \leanok
  Shift the chosen distinguished cofiber triangle, then use the triangle
  isomorphism with components $1,(-1)^n,1$ to make its first two arrows
  positive while retaining the signed connecting map. Contravariant
  representable exactness at the middle term gives the extension.
  No additional sign compatibility is assumed.
\end{proof}

\begin{proposition}[Cofiber comparison of an extension and its restriction]
  \label{prop:cofiber-extension-octahedron}
  \lean{KIP126.StableHomotopy.CofiberExtension.cofiber_triangle_of_extension,KIP126.StableHomotopy.CofiberExtension.exists_eta_nu_cofiber_triangles,KIP126.StableHomotopy.CofiberExtension.sphereSixIso,KIP126.StableHomotopy.CofiberExtension.shiftFourIso}
  \uses{prop:shifted-cofiber-extension,def:mathlib-octahedron}
  \leanok
  In the same category with the octahedron axiom, let
  $g:\Sigma^n C(f)\to Z$ restrict along the positive shifted inclusion
  to $a:\Sigma^nY\to Z$. There are maps
  $m:\Sigma^{n+1}X\to C(a)$ and $j:C(a)\to C(g)$ from one octahedron.
  They give the distinguished triangle
  \[
    C(a)\xrightarrow{j} C(g)
      \xrightarrow{(\Sigma\delta_{f,n})\delta_g}
      \Sigma^{n+2}X\xrightarrow{-\Sigma m}\Sigma C(a),
  \]
  where $\delta_{f,n}$ is the signed connecting map of the normalized
  shifted cofiber triangle. All four octahedron squares are retained,
  including $j\iota_a=\iota_g$ and
  $\delta_gj=(\Sigma^{n+1}\iota_f)\delta_a$.
  No independent choice of a cofiber map replaces $j$.
  For $f=\eta:S^1\to S^0$, $n=3$ and $a=\nu$, the explicit composite-zero
  premise supplies one $g,m,j$ and both triangles
  $S^0\to C(g)\to\Sigma^4C\eta$ and
  $C\nu\to C(g)\to S^6$. The former's second map is exactly
  $\delta_g$ followed by the canonical shift-composition isomorphism.
  Actual composite vanishing and native-module comparisons remain separate.
\end{proposition}
\begin{proof}
  \leanok
  Apply the octahedron axiom to the rotated normalized shifted triangle,
  the chosen cofiber triangle of $g$, and the chosen cofiber triangle of
  $a$. Rotate its new distinguished triangle and keep the same two maps
  and four squares. The final connecting arrow is $-\Sigma m$.
\end{proof}

% SOURCE: MainPaper/main.tex:253-271, 648-664, and 929-971; exactness used repeatedly.
\begin{theorem}[Stable homotopy long exact sequence]
  \label{thm:stable-homotopy-long-exact}
  \uses{def:stable-homotopy-group,def:cofiber-triangle,def:mathlib-homological-functor,thm:mathlib-homology-sequence-exact-one,thm:mathlib-homology-sequence-exact-two,thm:mathlib-homology-sequence-exact-three}
  \notready
  Every chosen cofiber triangle
  $X\to Y\to C(f)\to\Sigma X$ induces the natural long exact sequence
  \[
    \cdots\to\pi_{n+1}C(f)\to\pi_nX\to\pi_nY
      \to\pi_nC(f)\to\pi_{n-1}X\to\cdots,
  \]
  with connecting maps compatible with rotation, suspension, and morphisms of
  triangles.
\end{theorem}

\begin{proof}
  Equip the representable functor $[S^0,-]$ with Mathlib's homological-functor
  structure, apply its three homology-sequence exactness lemmas to every
  shift, and splice the rotated triangles.  Naturality follows from a morphism
  of triangles.
\end{proof}

% SOURCE: MainPaper/main.tex:1755-1921, smash diagrams in May's lemma and the Mahowald trick.
\begin{definition}[Smash of cofiber triangles]
  \label{def:smash-triangle}
  \uses{def:cofiber-triangle}
  \notready
  Smashing a distinguished triangle with any spectrum gives a distinguished
  triangle, naturally in both variables.  For two cofiber triangles the
  induced $3\times3$ diagram contains the canonical pullback/pushout square
  and connecting maps used to compare the two ways of reaching the double
  cofiber.
\end{definition}

\subsection{Homotopy classes and secondary operations}

% SOURCE: MainPaper/main.tex:2387-2616, Toda-bracket arguments.
\begin{definition}[Toda bracket]
  \label{def:toda-bracket}
  \uses{def:cofiber-triangle}
  \notready
  Given composable stable maps $a,b,c$ with chosen nullhomotopies of $ba$ and
  $cb$, the triple Toda bracket $\langle c,b,a\rangle$ is a coset of homotopy
  classes with its standard left and right indeterminacy.  The interface
  includes suspension, naturality, containment under changed nullhomotopies,
  and Moss-style comparison with Massey products when an Adams convergence
  theorem supplies that comparison.
\end{definition}

% SOURCE: MainPaper/main.tex:2470-2567, Massey products on Adams pages; Moss Theorem 1.2.
\begin{definition}[Massey product]
  \label{def:massey-product}
  \uses{def:multiplicative-ss,def:mathlib-f2}
  \notready
  On a differential graded page model for the multiplicative Adams spectral
  sequence, composable classes $a,b,c$ with $ab=0$ and $bc=0$ and chosen
  nullhomologies define a triple Massey-product coset
  $\langle a,b,c\rangle$.  The interface records its bidegree, left and right
  indeterminacy, naturality, change-of-choice formula, and a proposition that
  the indeterminacy is zero in a specified finite bidegree.
\end{definition}
